\section{The exact metric departure and the cost of complete jets}
\label{sec:metric-departure}
Retain the original packet \(E=\mathbb C[s]/h\), its polynomial
coordinates and \(A=M_s\). Write its eigenvalues with their algebraic
multiplicities as \(\rho_1,\ldots,\rho_d\), and put
\(\delta_j=\operatorname{Re}\rho_j-1/2\). No eigenvalue or repeated
factor is discarded. For a positive Hermitian matrix \(G\) in these
coordinates define
\[
 \begin{gathered}
 W_G=A^*G+GA-G,\quad S_G=G^{-1/2}W_GG^{-1/2},\quad
 \epsilon_G=\|S_G\|_{\mathrm{op}},\\
 \mathfrak d_G(A)=\operatorname{Tr}(G^{-1}A^*GA)
                          -\sum_{j=1}^d|\rho_j|^2 .
 \end{gathered}
 \tag{MD1}
\]
The positive square root is used only for the explicit isometry
\(x\mapsto G^{1/2}x\) from the specified metric space \((E,G)\)
to the Euclidean coordinate space. The original \(A,G\), their
entries and their maps remain those in (MD1).

\begin{theorem}[Exact departure identity]
The number \(\mathfrak d_G(A)\) is real and nonnegative, and
\[
 \operatorname{Tr}S_G=2\sum_j\delta_j,\qquad
 \operatorname{Tr}(S_G^2)=4\sum_j\delta_j^2+2\mathfrak d_G(A).
 \tag{MD2}
\]
It vanishes precisely when \(A\) is normal for the specified
metric, that is, when
\(A(G^{-1}A^*G)=(G^{-1}A^*G)A\).
For the original orthogonal representatives \(R_m\), their
matrices \(G_m\) and rank-two forms \(W_m\), one has exactly
\[
 \epsilon_m=\left|\sum_j\delta_j\right|
 +\sqrt{2\sum_j\delta_j^2+\mathfrak d_{G_m}(A)
                     -\left(\sum_j\delta_j\right)^2}.
 \tag{MD3}
\]
For a reflection-stable packet this becomes
\[
 \boxed{\epsilon_m^2=2\sum_{\rho\in Z}m_\rho
               (\operatorname{Re}\rho-1/2)^2+\mathfrak d_{G_m}(A).}
 \tag{MD4}
\]
\end{theorem}
\begin{proof}
Let \(C=G^{1/2}AG^{-1/2}\). Then
\(S_G=C+C^*-I\) and
\(\operatorname{Tr}(G^{-1}A^*GA)=\operatorname{Tr}(C^*C)\).
For completeness, choose a unit eigenvector of \(C\), extend it
to an orthonormal basis, and repeat on the lower-right compression.
Induction gives a unitary matrix \(U\) with upper-triangular
\(T=U^*CU\) whose diagonal lists the same eigenvalues with their
algebraic multiplicities. Unitary invariance of the trace gives
\[
 \mathfrak d_G(A)=\sum_{i<j}|T_{ij}|^2\geq0 .
 \tag{MD5}
\]
If it is zero, \(T\) is diagonal and \(C\) is normal. Conversely,
an upper-triangular normal matrix is diagonal: its first row norm
equals its first column norm, eliminating the strict first-row
entries, and induction eliminates every remaining such entry.
Thus normality of \(C\) is equivalent to the stated equality for
\(A\), through the specified isometry.

Expanding \((C+C^*-I)^2\) under the trace gives
\[
 \operatorname{Tr}(S_G^2)=2\operatorname{Tr}(C^*C)
                 +2\operatorname{Re}\operatorname{Tr}(C^2)
                 -4\operatorname{Re}\operatorname{Tr}C+d .
\]
The traces of \(C\) and \(C^2\) are respectively
\(\sum_j\rho_j\) and \(\sum_j\rho_j^2\), even in the presence
of nilpotents. Since
\(2|\rho|^2+2\operatorname{Re}\rho^2
-4\operatorname{Re}\rho+1=4(\operatorname{Re}\rho-1/2)^2\),
this proves (MD2). The trace formula follows directly from
\(S_G=C+C^*-I\).

For \(G_m\), Section~\ref{sec:arithmetic-input} proved
\(\operatorname{rank}S_{G_m}\leq2\). List its two possibly
nonzero real eigenvalues \(a,b\), appending zero if necessary
(also when \(d=1\)). Equations (MD2) give
\[
 a+b=2\sum_j\delta_j,\quad
 (a-b)^2/4=2\sum_j\delta_j^2+\mathfrak d_{G_m}(A)
                          -(\sum_j\delta_j)^2 .
\]
Thus the radicand is nonnegative, and
\(\max(|a|,|b|)=|(a+b)/2|+|a-b|/2\), proving (MD3).
Reflection-stability pairs \(\delta\) with \(-\delta\), with
the same multiplicities, proving (MD4).
\end{proof}

The actual arithmetic kernel of Section~\ref{sec:tau-boundary-kernel}
also gives this quantity without changing the local unit. Indeed
\(G_m=U_\varepsilon^*K_N^{-1}U_\varepsilon\), \(N=d+m\), and
\(U_\varepsilon A=AU_\varepsilon\). Cyclicity of trace gives
\[
 \mathfrak d_{G_m}(A)
   =\operatorname{Tr}(K_NA^*K_N^{-1}A)-\sum_j|\rho_j|^2.
 \tag{MD6}
\]
To verify the cancellation, substitute
\(G_m^{-1}=U_\varepsilon^{-1}K_N(U_\varepsilon^*)^{-1}\)
in (MD1), commute \(A\) with \(U_\varepsilon\) and \(A^*\)
with \(U_\varepsilon^*\), and cancel the outer similar factors
inside the trace. This calculation removes no source representative:
their exact recovery remains (TB.7)--(TB.9).

\subsection{A quantitative constraint on every retained nilpotent chain}
For a selected root \(\rho\) of full order \(r\), let
\[
 I_\rho:\mathbb C[z]/z^r\longrightarrow E
\]
be the original Chinese-remainder inclusion: the \(\rho\)-component
is the given Taylor polynomial in \(z=s-\rho\), and every other
component is zero. Its Gram matrix is
\(G_\rho=I_\rho^*GI_\rho\), in the basis
\(1,z,\ldots,z^{r-1}\). Put \(N_\rho z^k=z^{k+1}\), with
\(z^r=0\). The exact intertwining relation is
\(AI_\rho=I_\rho(\rho I+N_\rho)\).

\begin{theorem}[Full-jet metric inequalities]
For every positive original metric \(G\), write
\(\delta=\operatorname{Re}\rho-1/2\). Then
\(\epsilon_G\geq2|\delta|\), and for all
\(0\leq k<k+\ell<r\),
\[
 \frac{(G_\rho)_{k+\ell,k+\ell}}{(G_\rho)_{k,k}}
 \leq \left[\frac r2(\epsilon_G^2-4\delta^2)\right]^\ell.
 \tag{MD7}
\]
In particular, for \(r\geq2\),
\[
 \kappa_2(G_\rho)\ \geq\
 \left[\frac r2(\epsilon_G^2-4\delta^2)\right]^{-(r-1)} .
 \tag{MD8}
\]
The bracket is strictly positive when \(r\geq2\).
Here \(\kappa_2\) is the ratio of the largest and smallest
eigenvalues in the specified Taylor coordinates.
\end{theorem}
\begin{proof}
The definition of \(\epsilon_G\) is the two-sided form inequality
\(-\epsilon_GG\leq W_G\leq\epsilon_GG\).
On an eigenvector it gives \(\epsilon_G\geq2|\delta|\).
Pull this same inequality back by \(I_\rho\). Let
\(L=G_\rho^{1/2}N_\rho G_\rho^{-1/2}\) and
\(T=2\delta I+L+L^*\). The pulled-back inequality gives
\(\|T\|\leq\epsilon_G\). Nilpotence retains
\(\operatorname{Tr}L=\operatorname{Tr}(L^2)=0\), whence
\[
 \operatorname{Tr}T=2r\delta,\qquad
 \|L\|_F^2=\tfrac12\operatorname{Tr}(L+L^*)^2
       =\tfrac12\operatorname{Tr}T^2-2r\delta^2
       \leq\tfrac r2(\epsilon_G^2-4\delta^2).
 \tag{MD9}
\]
The first equality for the Frobenius norm follows by expanding
the square and using the two zero traces. The final inequality
uses the \(r\) real eigenvalues of \(T\), each bounded in modulus
by \(\epsilon_G\). For the original metric norm,
\(\|N_\rho\|_{G_\rho}=\|L\|\leq\|L\|_F\).
Apply its \(\ell\)-th power to the specified basis vector \(z^k\).
Since \(N_\rho^\ell z^k=z^{k+\ell}\), the resulting squared
norm inequality is exactly (MD7).

When \(r\geq2\), \(N_\rho\neq0\), so (MD9) makes the bracket
strictly positive. Put \(k=0,\ell=r-1\). The Rayleigh inequalities
give
\(\kappa_2(G_\rho)\geq (G_\rho)_{0,0}/
(G_\rho)_{r-1,r-1}\); combine with (MD7) to obtain (MD8).
\end{proof}

All entries in (MD7) for the original sequence are explicitly
\(I_\rho^*U_\varepsilon^*K_N^{-1}U_\varepsilon I_\rho\).
Thus this theorem supplies constraints on the actual moment
calculation, with its original Taylor factors. At a critical root
of order \(r\geq2\), a sequence with relative control approaching
zero necessarily has the stated diverging Taylor-metric condition
number. The surviving jets have not been erased; the inequalities
calculate their relative metric scales. This is compatible with
the complete graph and weighted-jet maps in
Section~\ref{sec:jet-topology}, which retain those same coordinates.
